\documentclass{cs22-resource}
\title{CS22 \LaTeX\ Guide}
\usepackage{listings}
\lstset{
    basicstyle=\sffamily,
    numbers=left,
    numberstyle=\tiny,
    frame=tb,
    columns=fullflexible,
    showstringspaces=false,
    language=[LaTeX]{TeX}
}

\begin{document}
\section*{What is LaTeX?}
\LaTeX\ is a typesetting system. Unlike most word processors, where what you
write is exactly what appears on your document, in \LaTeX\ you write your source
text separately (somewhat like HTML, or Markdown). Instead of manually formatting your document, in \LaTeX\ many
formatting details are standardized and taken care of automatically, and the
focus is placed on content. In addition, \LaTeX\ works more like code in that
there are many commands by which you can insert characters, format equations,
work with images and graphics, and much more. You can even define your own
commands to carry out more specialized functions, much like a programming
language.

\section*{Getting Started}
You can download \LaTeX\ software on to your computer, but we've found the
simplest way to start using \LaTeX\ is to use \href{https://www.overleaf.com}{Overleaf}. You can create
a free account with your Brown email (which includes additional features) and start making new projects right away.

Once you open a new project, you will see three areas on your screen. The
leftmost area is where you see your files; for the most part there is nothing to
do there unless you need to add images, in which case simply press the upload
button at the top and add in files from your computer. In the center you see
the source file, which is where you write your document in \LaTeX. Finally, on
the right is the final, compiled document, which is what your final product will
look like. Make sure you press the green \emph{Recompile} button to update the right
side periodically, and before you need to download your file.

When you are done writing the document, press the Menu button in the top
right corner to download your file(s). In general, for submissions you will only
need the PDF file representing your compiled work. The source file contains the
\LaTeX\ code you wrote in case you want to edit it elsewhere.

Here is a preview of the Overleaf editor layout and its components:

\includegraphics[width=1\textwidth]{overleaf_guide.jpg}

Of course, you may also look into installing \LaTeX\ locally on your computer and using a text editor like VSCode to edit and view your \LaTeX\ files. The department machines also all have \LaTeX\ installed (you can use FastX to connect to the department, and use an app like VSCode).

\section*{CS22 Document Class}
The TAs have created a \texttt{cs22} \LaTeX\ document class file so you can easily format your solutions. 

You should download the \href{https://gist.githubusercontent.com/jchen/69d56f5e45f868e165f09ec9059b8144/raw/b699e2e19b60d4e297843866641391798e19c90c/cs22.cls}{\texttt{cs22.cls}} file, and use it as a document class within your main \texttt{.tex} file. 

Your source code should look something like: 
\begin{lstlisting}
\documentclass{cs22}
\title{Homework 1}
\begin{document}
    \begin{problem}
        This is the solution to problem 1.
    \end{problem}
    \begin{problem}
        Likewise, this is the solution to problem 2.
    \end{problem}
\end{document}
\end{lstlisting}
When you compile your code, you should see a PDF appear with each problem solution on each page. This is formatted \emph{just like how we want you to submit your assignments!} Note how each solution is on a new page, which makes it easier to assign pages on Gradescope. 

This document class contains most packages you will need for CS22 (you can take a look for yourself). Nevertheless, you might want to add additional imports or commands to the \emph{preamble}. Between line 2 and line 3 (before you begin the document), you are able to add any package imports. More information about packages can be found \href{https://www.overleaf.com/learn/latex/Understanding_packages_and_class_files}{here}. 

One useful thing you can do is to define your own commands: 
\begin{center}
\texttt{\textbackslash newcommand\{[your shortcut]\}\{[actual command]\}}
\end{center}
This is essentially creating a shortcut for any commands you want to write. 

Your answers go between every \texttt{\textbackslash begin\{problem\}} and \texttt{\textbackslash end\{problem\}}. 

\section*{Learning LaTeX}
There are many resources online that will help you learn \LaTeX, and we think these are most appropriate for getting you started on your journey. The CS22 TAs will also be holding a \LaTeX\ gearup so you can get started with using the tool. 

Here are a selection of a few resources: 
\begin{itemize}
    \item \href{https://www.overleaf.com/learn/latex/Learn_LaTeX_in_30_minutes}{Learn LaTeX in 30 minutes - Overleaf}
    \item \href{https://www.cs.princeton.edu/courses/archive/spr10/cos433/Latex/latex-guide.pdf}{A Beginner’s Guide to LaTeX - David Xiao}
    \item \href{https://www.learnlatex.org/en/}{LearnLaTeX.org}
    \item \href{http://www.docs.is.ed.ac.uk/skills/documents/3722/3722-2014.pdf}{LaTeX for Beginners} (pretty longish book)
    \item \href{https://www.youtube.com/watch?v=VhmkLrOjLsw}{LaTeX Tutorial (YouTube) - Derek Banas} (hour-long video tutorial)
    \item \href{https://artofproblemsolving.com/wiki/index.php/LaTeX:Symbols}{LaTeX Symbols}
    \item \href{https://cs22.io/#resources}{More resources linked!}
    \item Google is always your friend!
\end{itemize}

\end{document}